% Abstract.
%
% Follows the rhythm of Hu et al. (CKA-RL, NeurIPS 2025) beat for beat, so it
% reads as native to the venue:
%   1. broad RL, not continual RL, one plain sentence;
%   2. "However," the non-stationarity that motivates the field;
%   3. "Although Continual Reinforcement Learning ..., existing methods ...";
%   4. "To address this, we propose NAME ...";
%   5. "Specifically, we introduce ..." the first component;
%   6. an effect sentence saying what that component buys;
%   7. "Additionally, ..." the second component;
%   8. "Experiments on three ... demonstrate ..." with concrete numbers;
%   9. code availability.
%
% Their abstract is 171 words, so the rhythm comes with an economy the first
% draft of this one did not have. This is 231, with no sentence over 33 words
% against their longest of 30.
%
% Two deliberate departures from their template. Beat 8 says "matches" rather
% than "outperforms state-of-the-art methods", since we are third of eleven on
% the established benchmark, and the clause naming what that benchmark
% suppresses is what earns the other two settings. The simulator assumption and
% the compute cost ride in beat 7 rather than waiting for a limitations
% section, per docs/storyline.md, because a reviewer who finds a
% stronger-than-replay assumption in an appendix reads it as concealment.
%
% Cut for length, and the first thing to restore if room appears: the Atari
% result, that DUEL alone both retains the earlier games and learns the last.
% "Three settings" currently carries it implicitly.

\begin{abstract}
Reinforcement learning enables agents to acquire effective behaviour through
interaction with an environment. However, real-world agents encounter tasks in
sequence, and learning a new one typically overwrites what earlier ones
established. Although Continual Reinforcement Learning addresses this setting,
existing methods only approximate retention, pricing it with a fixed penalty
weight, approximating it with a stored sample, or purchasing it with frozen
parameters. To address this, we propose DUEL, a dual-policy method that states
retention as an explicit constraint rather than a trade-off. Specifically, we
introduce a specialist shortfall constraint requiring the deployed policy to
stay within a tolerance of a local specialist trained on each task as it
arrives. The constraint is one-sided, so the policy cannot fall behind a
specialist yet may exceed one, which keeps positive backward transfer
available. Additionally, we enforce it by primal-dual alternation, so each
task's multiplier is set by the optimisation rather than in advance. DUEL
stores no trajectories, re-entering past-task simulators at roughly 1.5 times
the compute of the cheapest baseline. Experiments on three settings show that
DUEL matches state-of-the-art methods on the established benchmark, where
per-task heads and same-game tasks can substantially mask interference. Under
one shared head across fifty tasks it improves average performance by 0.126 and
ends 91\% of tasks above the level at which they were first learned, against
13\% for the nearest baseline. The source code will be released.
\end{abstract}
